Real-time optimisation (RTO) drives a process toward a quality constraint it cannot
measure online, so it must subtract a back-off, a safety margin, from the specification.
Conformal prediction is the natural distribution-free way to size that back-off from
data. We show that it is unsafe under optimisation. An RTO that optimises against a
marginally valid conformal back-off, even a locally adaptive one, behaves as a selection
mechanism: it drives the operating point to where the margin under-covers the conditional
constraint quantile, and on a rigorous Peng--Robinson distillation twin the realised
violation rate reaches roughly five times the nominal level across the disturbance range.
We identify this conformal selection effect as the optimisation-induced loss of marginal
coverage, and we verify against an oracle, which uses the true conditional quantile, that
the failure lies in the back-off rather than in the chance-constrained formulation.
Restoring safety requires conditional validity. A back-off built from conformalised
quantile regression, combined with an a-posteriori calibration step at the selected
setpoint, returns the realised violation to the oracle level, within a Monte-Carlo band,
at near-oracle profit over the operationally realistic disturbance range. It degrades only
when the disturbance widens several-fold and the conditional estimate runs out of
calibration data, and even then it degrades visibly, through a climbing inflation factor.
A regime map over disturbance magnitude bounds where each back-off controls violation. The
contribution is calibrated safety, not profit: at a well-controlled feed the constraint is
barely active and all methods earn within half a percent of the deterministic optimum. The
mechanism is general, since any conformal interval used as a hard optimisation constraint
is exposed to it, and the portable rule that follows is to size the margin conditionally
and to read the a-posteriori inflation as a signal of data adequacy.
